\documentclass[../../main.tex]{subfiles}% 注意这里的文件路径不能用 ./main.tex ,否则用latexmk编译子文件会报错
\graphicspath{{\subfix{./image/}}{\subfix{../../image/}}} % 指定图片引用路径,后续可以直接使用图片文件名
% 如果图片存放在上级目录,则使用 ../../image/ 作为备用路径

% 例如:
% \begin{figure}[H]
% \centering
% \includegraphics[width=0.3\textwidth]{图.png}
% \caption{}
% \label{figure:图}
% \end{figure}
% 注意:上述\label{}一定要放在\caption{}之后,否则引用图片序号会只会显示??.

\begin{document}

\section{适定性问题}

\begin{definition}
如果一个偏微分方程定解问题满足下列条件：
\begin{enumerate}[(1)]
\item 它的解存在；
\item 它的解惟一；
\item 它的解连续地依赖定解条件和定解问题中的已知函数，
\end{enumerate}
则称这个定解问题是\textbf{适定的}；否则称这个定解问题是\textbf{不适定的}。
\end{definition}

当求解位势方程、热方程和波动方程的定解问题时，我们总是先假设定解问题的解具有非常好的性质（如光滑性）；接着我们求出定解问题的解的表达式。这样的解我们称为定解问题的\textbf{形式解}。然后我们再严格证明在定解条件满足一定的要求时，所得到的形式解的确是定解问题的古典解，从而获得解的存在性。在研究位势方程、热方程和波动方程的定解问题的解的惟一性和稳定性时，我们先导出一些关于解的\textbf{先验估计}；然后再利用这些估计推导出定解问题的解的惟一性和稳定性。这样，我们就基本上回答了关于位势方程、热方程和波动方程的一些定解问题的适定性问题。

\begin{definition}
设 $\Omega\subset \mathbb{R}^n$,我们用 $\overline{\Omega}$ 表示它的闭包，$\partial\Omega$ 表示它的边界。同时，我们用
\begin{align*}
\mathbb{R}_+^n = \{x=(x_1,\cdots,x_n)\in\mathbb{R}^n \mid x_n > 0\}
\end{align*}
表示 $n$ 维欧氏空间 $\mathbb{R}^n$ 中的上半空间。由于上半空间 $\mathbb{R}_+^n$ 的边界
\begin{align*}
\partial\mathbb{R}_+^n = \{x=(x_1,\cdots,x_n)\in\mathbb{R}^n \mid x_n=0\}
\end{align*}
是 $n-1$ 维欧氏空间，我们也记 $\partial\mathbb{R}_+^n = \mathbb{R}^{n-1}$。当 $n=1$ 时，我们将 $\mathbb{R}_+^1$ 简记为 $\mathbb{R}_+$，因而 $\mathbb{R}_+^{n+1} = \mathbb{R}^n \times \mathbb{R}_+$。
\end{definition}

\begin{theorem}
\begin{enumerate}[(1)]
\item 我们用 $B(x,r)$ 表示 $n$ 维欧氏空间 $\mathbb{R}^n$ 中以 $x$ 为中心，以 $r$ 为半径的开球，其体积为
\begin{align*}
\alpha(n)r^n = \frac{2\pi^{\frac{n}{2}}}{n\Gamma\left(\frac{n}{2}\right)}r^n,
\end{align*}
其中 $\alpha(n)$ 表示 $\mathbb{R}^n$ 上单位球的体积.

\item 我们用 $S^{n-1}(x,r)=\partial B(x,r)$ 表示球 $B(x,r)$ 的边界，即 $n-1$ 维球面，其面积为
\begin{align*}
\omega(n)r^{n-1} = n\alpha(n)r^{n-1}.
\end{align*}
\end{enumerate}
\end{theorem}
\begin{proof}
\begin{enumerate}[(1)]
\item 只需计算单位球的体积.假设结论对小于$n$的情况都成立,现在考虑$n$的情形.
\begin{align*}
\alpha(n) &= |B(0,1)|
= \int_{-1}^{1} \left| B_{n-1}\left( (0,\cdots,0,z), \sqrt{1-z^2} \right) \right| \mathrm{d}z \\
&= 2\int_{0}^{1} \alpha(n-1)(1-z^2)^{\frac{n-1}{2}} \mathrm{d}z \xlongequal{u = z^2} 2\alpha(n-1)\int_{0}^{1} (1-u)^{\frac{n-1}{2}} \frac{1}{2} u^{-\frac{1}{2}} \mathrm{d}u \\
&= \alpha(n-1) B\left( \frac{n+1}{2}, \frac{1}{2} \right) = \alpha(n-1) \frac{\Gamma\left(\frac{n+1}{2}\right)\Gamma\left(\frac{1}{2}\right)}{\Gamma\left(\frac{n+2}{2}\right)}
\\
&=\alpha \left( n-2 \right) \frac{\Gamma \left( \frac{n}{2} \right) \Gamma \left( \frac{1}{2} \right)}{\Gamma \left( \frac{n+1}{2} \right)}\frac{\Gamma \left( \frac{n+1}{2} \right) \Gamma \left( \frac{1}{2} \right)}{\Gamma \left( \frac{n+2}{2} \right)}=\alpha \left( 1 \right) \prod_{k=2}^{n}{\frac{\Gamma \left( \frac{k+1}{2} \right) \Gamma \left( \frac{1}{2} \right)}{\Gamma \left( \frac{k+2}{2} \right)}}
\\
&= \frac{\Gamma\left(\frac{3}{2}\right)\left[\Gamma\left(\frac{1}{2}\right)\right]^{n-1}}{\Gamma\left(\frac{n+2}{2}\right)} \cdot \frac{2\pi^{\frac{1}{2}}}{\Gamma(\frac{1}{2})} = \frac{\frac{1}{2}\left[\Gamma\left(\frac{1}{2}\right)\right]^{n}}{\Gamma\left(\frac{n}{2} + 1\right)} \cdot 2 \quad (\Gamma(s+1) = s\Gamma(s), s > 0) \\
&= \frac{\pi^{\frac{n}{2}}}{\Gamma\left(\frac{n}{2} + 1\right)}=\frac{2\pi^{\frac{n}{2}}}{n\Gamma(\frac{n}{2})} \quad \left( \Gamma(s)\Gamma(1-s) = \frac{\pi}{\sin \pi s}, 0 < s < 1 \right).
\end{align*}

\item 只需计算单位球的面积.利用变量替换$(x_1,x_2,\cdots,x_n)\mapsto (ry_1,ry_2,\cdots,rt_{n-1},x_n)$可得
\begin{align*}
\alpha(n) = \int_{0}^{1} \omega(n) r^{n-1} \mathrm{d}r = \frac{\omega(n)}{n},
\end{align*}
整理即得 $\omega(n) = n\alpha(n)$.
\end{enumerate}
\end{proof}




\end{document}